\vspace{+0.5em}
\section{Education}
%\cventry{2023--Current}{PhD in Mechanics and Advanced Engineering Sciences}{University of Bologna (DIMSAI)}{}{\textit{}Leveraging IoT-generated Multimodal Context-Aware system to promote Perceptual-Motor Skills in Development Coordination Disorders (DCD)}{}  % Arguments not required can be left empty
%\eduentry{2023--Mar 27}
%{PhD, Mechanics and Advanced Engineering Sciences}
%{\href{https://www.unibo.it/}{\color{blue}Bologna University}}
%{Italy}
%{Topic: IoT-based Multimodal Context-Aware System – OmniTrack}
%Job Summary:
%We are seeking a talented professional to join our team as an IoT Solutions Architect / Connected Devices Engineer / Sensor Data Analyst. The ideal candidate will design, implement, and analyze IoT systems, ensuring seamless integration of connected devices and actionable insights from sensor data.Key Responsibilities:
%IoT Solutions Architecture:Design end-to-end IoT systems and architectures to meet business requirements.
%Select appropriate hardware, connectivity, and cloud platforms for IoT deployments.
%Ensure security, scalability, and reliability of IoT solutions.
%Connected Devices Engineering:Develop, configure, and maintain IoT devices and sensors.
%Implement firmware, connectivity protocols, and device integration solutions.
%Troubleshoot and optimize device performance and reliability.
%Sensor Data Analysis:Collect, process, and analyze data from connected devices.
%Develop actionable insights and reports to support business decision-making.
%Collaborate with cross-functional teams to implement data-driven solutions.
%Qualifications:
%Bachelor’s degree in Computer Science, Electrical Engineering, IoT, or related field.
%Proven experience in IoT architecture, device engineering, or sensor data analysis.
%Hands-on experience with IoT platforms, cloud integration (AWS IoT, Azure IoT, GCP IoT).
%Proficiency in programming languages such as Python, C/C++, or JavaScript.
%Strong analytical and problem-solving skills.
%Excellent communication skills to work with technical teams and business stakeholders.
%Preferred Skills: Experience with industrial IoT or smart building solutions.
%Knowledge of data analytics, machine learning, or AI applied to IoT data.
%Familiarity with connectivity protocols like MQTT, LoRaWAN, or Zigbee.
%Work Environment & Benefits:
%Flexible work arrangements: part-time, full-time, or hybrid.
%Opportunity to work with cutting-edge IoT and data analytics technologies.
%Professional growth and learning opportunities in IoT and connected devices.
%Collaborative and innovative team culture.

%\eduentry{April--Jul 26}
%{Mobility, Experimental Psychology}
%{\href{https://www.uu.nl/en}{\color{blue}Utrecht University}}
%{Utrecht, Netherlands}
%{Topic: Sensory navigation profile, a novel tool to predict successful use of sensory information }
%\eduentry{Nov--Dec 25}
%{Mobility, Information networking for innovation}
%{\href{https://www.iniad.org/en/}{\color{blue}Toyo University}}
%{Tokyo, Japan}
%{Topic: User-Centric Feedback Adaptation System (UCFAS): Analyzing Feedback Influences on User Behaviour in Navigation Tasks}
\eduentry{2023--27}
{PhD, Mechanics and Advanced Engineering Sciences}
{\href{https://www.unibo.it/}{\color{blue}Bologna University}}
{IT}
{Topic: IoT-based Multimodal Navigation System – OmniTrack\newline
• Research Mobility, \href{https://www.uu.nl/en}{\color{blue}Utrecht University}, NLD: Sensory navigation and experimental psychology.\newline
• Research Mobility, \href{https://www.iniad.org/en/}{\color{blue}Toyo University}, JP: Feedback algorithm and behaviour analysis.
}
\eduentry{2017--20}
{MSc, Desig\&Engineer}
{\href{https://www.polimi.it/en}{\color{blue}Politecnico di Milano}}
{IT}
{{Topic: UX development of an ultrasonic dental device: from concept to validation: WashB}\newline
• Research Mobility, \href{https://www.chiba-u.ac.jp/e/}{\color{blue}Chiba University}, JP: UX research studies for Fujitsu and Panasonic.
} 
%\eduentry{Mar--Sep 19}
%{Erasmus+}
%{\href{https://www.chiba-u.ac.jp/e/}{\color{blue}Chiba University}}
%{Japan}
%{Activity: Human-centred design approach in collaboration with Fujitsu and Panasonic.}

\eduentry{2012--15}
{BSc, Industrial Design}
{\href{https://www.unifi.it/changelang-eng.html}{\color{blue}Florence University}}
{IT}
{Topics: Foundation in research methodologies, development, and critical thinking.}
%\eduentry{2007--12}
%{HS, Computer Science}
%{\href{https://www.luigidellerba.edu.it}{\color{blue}ITIS Luigi dell'Erba}}
%{Bari, Italy}
%{Subjects: Mathematics, Statistics, Programming Languages (C, C++, Java, JavaScript)}



